\chapter{Коэрцитивность высокоуглового хвоста вихря}
\label{ch:v6-bosonic-defect-polar-high-angular-coercivity-gate}

\section{Проблема}

Полярный расчёт проверил только конечное окно $|m|\le8$. Рост нескольких
последних уровней ещё не доказывает устойчивость всех угловых чисел. Требуется
соединить прямой численный мост с аналитической нижней оценкой бесконечного
хвоста.

\section{Численный мост}

Тем же полным пятикомпонентным полярным гессианом на сетке $120$ узлов
проверены все целые числа
\[
  9\le |m|\le86.
\]
Минимальный уровень достигается при $|m|=9$ и равен
\[
  \lambda_9=5.1084779443.
\]
Затем последовательность строго возрастает и при $|m|=86$ достигает
\[
  \lambda_{86}=26.8490162062.
\]
Отрицательных уровней в мосте нет.

\section{Аналитический хвост}

Для остатка удержана половина положительного квадрата Ходжа связности, а его
сцепление с фазовой компонентой ограничено неравенством Юнга. Радиальные
квадраты производных отброшены как неотрицательные, но фоновые перекрёстные
члены полной второй вариации сохранены. После нормировки локальной кинетической
метрикой получена эрмитова матрица $L_m(r)$.

Для каждой строки применена оценка Гершгорина. Коэффициенты матрицы не выше
второй степени по $|m|$, поэтому нижняя граница имеет вид
\[
  g_{r,j}(m)=a_{r,j}m^2+b_{r,j}m+c_{r,j}.
\]
На $7999$ радиальных точках найдено
\[
 \min a_{r,j}=0.00125>0,
\]
а при $|m|=87$
\[
 \min g_{r,j}(87)=0.2113888889,
 \qquad
 \min g'_{r,j}(87)=0.0816666667.
\]
Следовательно, оценка положительна и возрастает для каждого целого
$|m|\ge87$. Частично-волновое разложение калибровочно-скалярного оператора и
центробежный рост высоких каналов согласуются со стандартным анализом вихря
Нильсена--Олесена \cite{high-angular-baacke-kevlishvili}.

\begin{theorem}
Эффективный пятикомпонентный гессиан исправленного вихря не имеет внутренних
отрицательных мод ни при одном целом угловом числе. После удаления переносной
пары $m=\pm1$ глобальный внутренний минимум равен
$4.1446985802$ и достигается в осесимметричном канале.
\end{theorem}

\section{Нулевая гипотеза и стоп-критерий}

Эффективная амплитудно-калибровочная подсистема теперь закрыта по всем угловым
числам. Нулевая гипотеза следующего гейта: дополнительные компоненты спинов
два и три не создают отрицательного полярного блока после использования уже
выведенных жёсткостей $Z_Q=Z_T=1/3$.

Стоп-критерий: любая сходящаяся отрицательная мода полного тензорного оператора
закрывает прямую нить. Если полный оператор нельзя однозначно собрать из
родительского действия, результат эффективной подсистемы не повышается до
сертификата всей модели.

\section{Статус}

Угловая устойчивость эффективного пятикомпонентного ядра закрыта для всех
$m\in\mathbb Z$. Полный гессиан $Q+T+B$, продольные тензорные каналы,
замкнутая петля и рождение наблюдаемой материи остаются открытыми.

\begin{thebibliography}{9}
\bibitem{high-angular-baacke-kevlishvili}
J.~Baacke, N.~Kevlishvili,
\emph{One-loop corrections to the string tension of the vortex in the
Abelian Higgs model}, arXiv:0806.4349.
\end{thebibliography}